\section{Mapa de la situación de aprendizaje}

{\fontsize{22}{26}\selectfont\bfseries\color{titular}Situación de aprendizaje\par}
\vspace{6pt}

% ---------- 1 y 2 ----------
\begin{filatarjetas}
  \begin{tarjeta}{1 Identificación}
    \campo{Título:} \tituloSA

    \campo{Curso y materia:} 1.º de Bachillerato (Ciencias y Tecnología), Matemáticas I.

    \campo{N.º de sesiones / temporalización:} 20 sesiones de 55 minutos en 5 semanas (4 sesiones semanales). Es evaluable.

    \campo{Autor/a:} \autor

    \campo{Mapa (idea central):} con vectores, trigonometría y rectas podemos medir lo que no alcanzamos y decidir con datos dónde generar energía solar en el propio instituto.
  \end{tarjeta}
  \begin{tarjeta}{2 Contexto y reto}
    \campo{Situación / contexto:} el IES Ejemplo (centro ficticio de la Comunidad de Madrid) tiene paneles solares en la cubierta del gimnasio y quiere ampliar la instalación. Antes necesita saber qué sombras proyectan sus edificios, pero nadie ha medido las alturas y no se puede subir a la cubierta.

    \campo{Pregunta-reto:} ¿cómo podemos medir lo que no podemos alcanzar y usar esas medidas para decidir dónde caben más paneles sin que se hagan sombra?

    \campo{Lo que les cuento el primer día:} «La dirección os encarga un informe de verdad. Hace más de dos mil años, Eratóstenes midió la Tierra con una sombra; vosotros vais a medir el instituto con un clinómetro y a defender vuestra propuesta ante quien decide».
  \end{tarjeta}
\end{filatarjetas}

% ---------- 3 y 4 ----------
\begin{filatarjetas}
  \begin{tarjeta}[Matriz curricular 1]{3 Saberes básicos / criterios de evaluación}
    \begin{tablaSA}{colspec={X[1.5,l]X[1,l]}}
      Saberes básicos (bloque) & Criterios de evaluación \\
      Vectores: operaciones, módulo, ángulo, proyección (A) & 7.1, 2.1, 5.2 \\
      Razones, identidades y ecuaciones trigonométricas (B y D) & 3.2, 3.1, 1.2 \\
      Resolución de triángulos, teoremas del seno y del coseno (B) & 1.1, 1.2, 2.1 \\
      Ecuaciones de la recta, posición, ángulos y distancias (C) & 5.1, 5.2, 7.2 \\
      Modelización y herramientas digitales (C y D) & 1.1, 2.2, 7.2, 8.1, 6.1 \\
      Actitudes y aportación de las matemáticas (F) & 9.1, 9.2, 9.3, 6.2 \\
    \end{tablaSA}

    \pista{Texto literal y operativización completa en la sección~\ref{sec:matrices}.}
  \end{tarjeta}
  \begin{tarjeta}[Matriz curricular 2]{4 Competencias específicas y criterios}
    \begin{tablaSA}{colspec={Q[l,wd=30pt]X[l]}}
      CE & Criterios de evaluación \\
      CE1 & 1.1, 1.2: modelizar y obtener soluciones \\
      CE2 & 2.1, 2.2: verificar y elegir la solución \\
      CE3 & 3.1, 3.2: conjeturar con tecnología \\
      CE5 & 5.1, 5.2: conectar ideas matemáticas \\
      CE6 & 6.1, 6.2: conectar con otras áreas y retos \\
      CE7 & 7.1, 7.2: representar con tecnologías \\
      CE8 & 8.1: comunicar con rigor \\
      CE9 & 9.1, 9.2, 9.3: error, perseverancia, equipo \\
    \end{tablaSA}
  \end{tarjeta}
\end{filatarjetas}

% ---------- 5 y 6 ----------
\begin{filatarjetas}
  \begin{tarjeta}[Matriz curricular 3 · descriptores operativos]{5 Perfil de salida}
    \begin{tablaSA}{colspec={Q[l,wd=24pt]Q[l,wd=44pt]X[l]}}
      CE & Descriptor & Cómo se moviliza \\
      CE1 & STEM3 & Diseñar en equipo el modelo de sombras \\
      CE2 & STEM2 & Repetir medidas y valorar el error \\
      CE3 & STEM1 & De observar casos a demostrar \\
      CE5 & STEM1 & Resolver por varias vías conectadas \\
      CE6 & CC4 & Relacionar propuesta y sostenibilidad \\
      CE7 & CD2 & Crear el libro respetando la autoría \\
      CE8 & CCL1 & Defender la propuesta ante la dirección \\
      CE9 & CPSAA3.2 & Reparto equitativo de roles \\
    \end{tablaSA}
  \end{tarjeta}
  \begin{tarjeta}{6 Producto final}
    \campo{Producto:} libro de GeoGebra por equipo con la medición real, el modelo de sombras y una propuesta razonada de ampliación de paneles.

    \campo{Formato y destinatario:} libro digital y defensa oral de unos cinco minutos ante la clase y la dirección del centro, que recibe las propuestas. Los datos se comparten con el centro socio de eTwinning.

    \campo{Agrupamiento:} equipos heterogéneos de 4 estudiantes con roles rotatorios\soloProducto{ (sección~\ref{sec:cooperativo})}; trabajo individual en fichas, cuestionarios y examen.

    \campo{ODS:} 7 (energía asequible y no contaminante) y 11 (ciudades y comunidades sostenibles).
  \end{tarjeta}
\end{filatarjetas}

% ---------- 7 ----------
\apartado{7 Secuencia de actividades de aprendizaje}

\begin{tablaSAlarga}{
  colspec={Q[l,wd=44pt]X[2.2,l]X[1.1,l]X[1.3,l]X[1.4,l]},
  rowsep=3pt, cells={font=\scriptsize},
  column{1}={bg=rosa, fg=urjc, font=\scriptsize\bfseries},
  row{1}={bg=urjc, fg=white, font=\scriptsize\bfseries},
}
  Fase & Actividades / tareas & Sesiones y agrupamiento & Recursos & {Evaluación\\(instrumento y criterio)} \\
  Inicio & Presentación del reto y activación de conocimientos previos (triángulo rectángulo, Pitágoras y coordenadas). & S1. Gran grupo e individual & Plano del centro, GeoGebra Classroom & Cuestionario diagnóstico (sin calificación) \\
  Desarrollo & \textbf{Ejercicios:} tras cada exposición (vectores, trigonometría, triángulos, rectas), práctica guiada y una pieza del reto aplicada al propio centro. & S3, S4 y S6--S11. Magistral, individual y equipos & Pizarra, GeoGebra, hoja de ejercicios, plano del centro & Cuestionarios Q1--Q4 (7.1, 5.2, 1.2, 2.1) \\
  & \textbf{Actividades:} talleres de GeoGebra para explorar, conjeturar y aplicar: vectores, goniometría, y rayo y sombra. & S2, S5, S12 y S13 (F3, en S9). Individual y equipo & GeoGebra, fichas F1--F4, tablero Kanban & Fichas (7.1, 2.1, 3.1, 3.2, 1.1, 1.2, 7.2); mapa conceptual (5.1) \\
  & \textbf{Tareas:} medir el centro, construir el modelo de sombras y elaborar la propuesta de ampliación. & S14 a S17. Equipo & Clinómetro, cinta métrica, GeoGebra, PVGIS, eTwinning & Libro y propuesta (1.1, 2.2, 7.2, 6.1, 6.2); observación (9.1, 9.3) \\
  Cierre & Examen individual; defensa ante la dirección, coevaluación y metacognición. & S18 individual; S19 y S20 equipos & Examen modelo, libros, rúbricas & Examen; exposición (8.1, 1.1, 2.2, 6.1, 6.2); co/autoevaluación (7.2, 8.1, 9.x) \\
\end{tablaSAlarga}

% ---------- 8 ----------
\apartado{8 Evaluación}

\begin{tablaSAlarga}{
  colspec={X[1.2,l]X[2,l]X[1.3,l]}, cells={font=\scriptsize},
  row{1}={bg=rosa, fg=urjc, font=\scriptsize\bfseries},
}
  Criterios de evaluación & Instrumentos (rúbrica, lista de cotejo, diario…) & {Agente y momento\\(auto/co/hetero)} \\
  1.1, 1.2, 2.1, 3.1, 5.1, 5.2, 7.1, 7.2 & Examen escrito, cuestionarios y fichas GeoGebra & Heteroevaluación, durante los talleres y en S18 \\
  1.1, 2.2, 6.1, 6.2, 7.2, 8.1 & Rúbricas del libro y propuesta y de la exposición & Hetero y coevaluación, S16, S19 y S20 \\
  3.2, 5.1, 9.1, 9.2, 9.3 & Lista de cotejo, mapa conceptual y co/autoevaluación & Hetero, co y autoevaluación, S2--S20 \\
\end{tablaSAlarga}

\par\vspace{10pt}

% ---------- 9 ----------
\begin{tarjeta}{9 Atención a la diversidad (DUA)}
  Varias formas de acceder al contenido (explicación, construcción dinámica y material escrito), de expresarlo (libro, oral, escrito) y de implicarse (roles y elección del objeto que se mide). Refuerzo por criterio en S13 y ampliación con el radio de la Tierra. Detalle en la sección~\ref{sec:diversidad}.
\end{tarjeta}
